% unidades/11-pasado.tex
\chapter{⌛ 過去形 — Pasado}
\unidadheader{11}{過去形}{Pasado}{Hablar de acciones y estados terminados}

\begin{tcolorbox}[vocabbox, title={📌 Vocabulario Nuevo}]
\begin{tabularx}{\linewidth}{llX}
\toprule
\textbf{Japonés} & \textbf{Español} & \textbf{Tipo} \\
\midrule
\vocabitem{昨日}{きのう}{ayer}{adv.}
\vocabitem{先週}{せんしゅう}{la semana pasada}{adv.}
\vocabitem{先月}{せんげつ}{el mes pasado}{adv.}
\vocabitem{去年}{きょねん}{el año pasado}{adv.}
\vocabitem{〜ました}{—}{pasado cortés (afirmativo)}{v. term.}
\vocabitem{〜ませんでした}{—}{pasado cortés (negativo)}{v. term.}
\vocabitem{〜かったです}{—}{pasado de adj-い}{adj. term.}
\vocabitem{〜でした}{—}{pasado de sust. / adj-な}{term.}
\bottomrule
\end{tabularx}
\end{tcolorbox}

\subsection*{🗣️ Frases Clave}

\frase{きのう\ruby{晩}{ばん}ごはんを\ruby{食}{た}べませんでした。}{Ayer no cené.}
\frase{\ruby{先週}{せんしゅう}\ruby{日本}{にほん}へ\ruby{行}{い}きました。}{Fui a Japón la semana pasada.}
\frase{テストはとても\ruby{難}{むずか}しかったです。}{El examen estuvo muy difícil.}
\frase{\ruby{昨日}{きのう}は\ruby{雨}{あめ}でした。}{Ayer llovió (fue lluvia).}

\subsection*{⚙️ Gramática}

\patron{Pasado de verbos}
  {Afirmativo: 〜ました \quad | \quad Negativo: 〜ませんでした}
  {\ej{\ruby{勉}{べん}強しました。}{Estudié.}
  \ej{\ruby{飲}{の}みませんでした。}{No bebí.}}

\patron{Pasado de adjetivos y sustantivos}
  {Adj-い: quitar い → かったです \quad | \quad Sust. / Adj-な: 〜でした}
  {\ej{\ruby{良}{よ}かったです。}{Estuvo bien.}
  \ej{きれいでした。}{Estaba hermoso.}}

\begin{tcolorbox}[ejerciciobox, title={📝 Ejercicios Rápidos}]
\begin{enumerate}
  \item Pasa a pasado: \ruby{行}{い}きます.
  \item Pasa a pasado: \ruby{高}{た}いです.
  \item Traduce: \ruby{去年}{きょねん}メキシコへ\ruby{来}{き}ました。
\end{enumerate}
\vspace{0.4em}
\textbf{Respuestas:}
1) \ruby{行}{い}きました\quad
2) \ruby{高}{た}かったです\quad
3) El año pasado vine a México.
\end{tcolorbox}

\begin{tcolorbox}[kanjibox, title={🀄 Kanjis de la Unidad}]
\begin{tabularx}{\linewidth}{cllXX}
\toprule
\textbf{Kanji} & \textbf{On} & \textbf{Kun} & \textbf{Significado} & \textbf{Vocab} \\
\midrule
\kanjirow{見}{けん}{み}{ver}{\ruby{見}{み}る / \ruby{見学}{けんがく}}
\kanjirow{聞}{ぶん・もん}{き}{escuchar / preguntar}{\ruby{聞}{き}く / \ruby{新聞}{しんぶん}}
\kanjirow{買}{ばい}{か}{comprar}{\ruby{買}{か}う / \ruby{買}{か}い\ruby{物}{もの}}
\kanjirow{読}{どく}{よ}{leer}{\ruby{読}{よ}む / \ruby{読書}{どくしょ}}
\kanjirow{書}{しょ}{か}{escribir}{\ruby{書}{か}く / \ruby{辞書}{じしょ}}
\bottomrule
\end{tabularx}
\end{tcolorbox}
